\documentclass[11pt]{article}
\usepackage[margin=1in]{geometry}
\usepackage{xcolor}
\usepackage{tcolorbox}
\usepackage{hyperref}
\usepackage{enumitem}
\usepackage{listings}

% Define colors
\definecolor{osaorange}{RGB}{220,115,38}
\definecolor{aiblue}{RGB}{0,83,155}

\title{\textbf{Group 4: Vitamin D Supplementation Study}\\
\large{Dataset Guidance}}
\author{}
\date{}

\begin{document}

\maketitle

\section*{Dataset Overview}

\textbf{File:} \texttt{vitamin\_d.csv}

\textbf{Sample Size:} n = 45 adults (15 per dose group)

\textbf{Research Context:} Randomized controlled trial comparing three vitamin D supplement doses (Low, Medium, High) on vitamin D blood levels after 8 weeks.

\subsection*{Variables}

\begin{itemize}
\item \texttt{id}: Participant identifier (1-45)
\item \texttt{group}: Supplement dose group (Low, Medium, High)
\item \texttt{baseline}: Baseline vitamin D level before supplementation (ng/mL)
\item \texttt{vitamin\_d}: Vitamin D level after 8 weeks of supplementation (ng/mL)
\item \texttt{age}: Age in years
\end{itemize}

\subsection*{Loading the Data in R}

\begin{lstlisting}[language=R, basicstyle=\small\ttfamily, breaklines=true]
# Load the dataset
data <- read.csv("vitamin_d.csv")

# View structure
str(data)
head(data)

# Summary by group
tapply(data$vitamin_d, data$group, summary)
\end{lstlisting}

\section*{Research Question to Consider}

\textbf{``Do different vitamin D supplement doses lead to different vitamin D blood levels?''}

Alternatively: ``Which supplement dose is most effective?''

\section*{Key Decision Point}

\begin{tcolorbox}[colback=osaorange!10,colframe=osaorange,boxrule=1pt,arc=2pt]
\textbf{The Challenge:} You have baseline AND final measurements for each person. How will you use the baseline information in your analysis using methods covered in this course?

\textbf{Questions to explore with AI:}
\begin{itemize}
\item What are different ways to handle baseline measurements?
\item Should you analyze final levels? Change scores? Both?
\item What methods from this course can compare three groups?
\end{itemize}

\textbf{Important note:} AI may suggest approaches for handling baseline that you haven't learned yet. Consider what's possible with methods covered in this course.
\end{tcolorbox}

\section*{Things to Think About}

\textbf{Exploratory Analysis:}
\begin{itemize}
\item Create visualizations (boxplots by group for final levels, baseline, or change)
\item Calculate summary statistics by supplement dose group
\item Consider: What question are you trying to answer about the doses?
\end{itemize}

\textbf{Statistical Methods:}
\begin{itemize}
\item Review methods for comparing groups
\item Consider: How many groups are you comparing?
\item Consider: What outcome should you analyze?
\end{itemize}

\textbf{Baseline Measurements:}
\begin{itemize}
\item You have baseline vitamin D levels - how could you use this information?
\item Different approaches: ignore baseline, calculate change, adjust for baseline
\item Which approaches are possible with methods you've learned?
\end{itemize}

\section*{AI Consultation Strategy}

\begin{tcolorbox}[colback=aiblue!10,colframe=aiblue,boxrule=1pt,arc=2pt]
\textbf{Important:} When you ask AI about studies with baseline measurements, it may suggest advanced methods beyond this course.

\textbf{Prompts to try:}
\begin{itemize}
\item ``How should I analyze this vitamin D supplement study?''
\item ``What's the best way to handle baseline measurements in this analysis?''
\item ``I haven't learned [method AI suggested]. Can I answer my question with ANOVA?''
\item ``What's the difference between analyzing final levels vs. change scores?''
\end{itemize}

\textbf{Watch for:}
\begin{itemize}
\item Tools suggesting different approaches to baseline measurements
\item Methods you haven't encountered in this course
\item Trade-offs between simple and complex approaches
\end{itemize}

\textbf{Document everything:} Record what each AI tool recommends - this is material for your AI comparison section!
\end{tcolorbox}

\section*{Questions for Your Analysis}

As you work through this dataset, consider:

\begin{enumerate}
\item What research question are you answering?
\item Which statistical method(s) from this course can address this question?
\item How are you handling the baseline measurements? Why?
\item What assumptions do you need to check?
\item If AI suggested advanced methods for baseline adjustment, why did you choose your approach?
\end{enumerate}

\section*{Possible Approaches to Explore}

Think about these options (all use methods from this course):

\begin{itemize}
\item \textbf{Option 1:} Analyze final vitamin D levels only
\item \textbf{Option 2:} Calculate change scores and analyze those
\item \textbf{Option 3:} Compare both approaches

Which makes most sense for your research question? Why?
\end{itemize}

\section*{Resources}

\begin{itemize}
\item Week 9 drop-in session (Wed 12-1:20pm, 153 Milam)
\item Textbook: Chapter 12 (Analysis of Variance)
\item Your team!
\item Canvas discussion board
\end{itemize}

\vspace{0.5cm}

\begin{tcolorbox}[colback=gray!10,colframe=gray!50,boxrule=1pt,arc=2pt]
\textbf{Remember:} Multiple valid approaches may exist. The goal is to:
\begin{itemize}
\item Choose a method you understand from this course
\item Justify why your approach answers the research question
\item Explain your reasoning about baseline measurements
\item Critically evaluate what different AI tools suggested
\item Document your decision-making process in your report
\end{itemize}
\end{tcolorbox}

\end{document}
